\section{Scalability Experiments \& Results}
\label{sec:scalability}

\subsection{Experimental Design}
\label{subsec:experimental-design}
The objective of our experiments was to evaluate the performance characteristics and limitations of the Spark-based data pipeline. We manipulated three independent variables across eight distinct trials: cluster size (1 to 3 workers), compute density (1 to 2 cores per executor), and workload volume (5~GB and 10~GB). This structured approach allowed us to isolate three specific scaling paradigms essential for distributed systems analysis:

\begin{itemize}
    \item \textbf{Horizontal (Strong) Scaling:} This measures the system's ability to reduce total execution time as additional worker nodes are added (Trials H-1, H-2, and H-3) while maintaining a fixed 5~GB workload. 
    \item \textbf{Vertical Scaling:} This evaluates the impact of increasing CPU cores per worker (1 vs. 2 cores) while keeping the cluster size fixed at three nodes. This was tested for both 5~GB (V-1, V-2) and 10~GB (V-3, V-4) datasets to identify if the job is compute-bound.
    \item \textbf{Weak Scaling:} This measures how the system handles proportional growth in both workload and resources (W-1 vs. W-2). Ideally, if the data per worker remains constant (5~GB/node), the runtime should remain flat.
\end{itemize}

The primary metrics for evaluation are total runtime in seconds, Speedup ($S = T_1/T_N$), and Efficiency ($E = S/N$).

\subsection{Horizontal Scaling (Strong Scaling)}
Using a fixed 5~GB workload and 2 cores per worker, we observed a clear trend of diminishing returns as the cluster size increased from one to three nodes:

\begin{itemize}
    \item \textbf{1 Worker (Trial H-1):} 142 seconds.
    \item \textbf{2 Workers (Trial H-2):} 89 seconds (A significant 53-second reduction).
    \item \textbf{3 Workers (Trial H-3):} 81 seconds (A marginal 8-second reduction).
\end{itemize}

\begin{figure}[H]
    \centering
    \includegraphics[width=0.65\textwidth]{horizontal_scaling_bar.png}
    \caption{Strong Scaling: Execution time for 5~GB dataset across 1, 2, and 3 workers.}
    \label{fig:horizontal_scaling}
\end{figure}

\textbf{Analysis:} The results indicate that while the jump from one to two workers yielded $\mathtt{\sim}$37\% improvement in speed, the transition to a third worker provided only $\mathtt{\sim}$9\% gain. This non-linear behavior suggests that for a relatively small 5~GB dataset, the administrative overhead of the Spark driver, specifically task distribution and managing RDD partitions, begins to consume a larger percentage of the total runtime than the actual computation.

\subsection{Vertical Scaling}
We compared 1 vs. 2 cores on a fixed 3-worker cluster to determine the impact of compute density on different data scales:

\begin{itemize}
    \item \textbf{5~GB Workload (V-1 vs. V-2):} Runtime decreased from 83s to 80s, a negligible gain.
    \item \textbf{10~GB Workload (V-3 vs. V-4):} Runtime decreased from 102s to 80s, a significant 21\% improvement.
\end{itemize}

\begin{figure}[H]
    \centering
    \begin{minipage}{0.48\textwidth}
        \centering
        \includegraphics[width=\textwidth]{vertical_scaling_5gb.png}
        \caption{Vertical Scaling (5~GB)}
        \label{fig:vertical_5gb}
    \end{minipage}
    \hfill
    \begin{minipage}{0.48\textwidth}
        \centering
        \includegraphics[width=\textwidth]{vertical_scaling_10gb.png}
        \caption{Vertical Scaling (10~GB)}
        \label{fig:vertical_10gb}
    \end{minipage}
\end{figure}

\textbf{Analysis:} The effectiveness of vertical scaling is strictly dependent on the data volume. At 5~GB, the workers are not CPU-bound; hence, additional cores do not help. However, at 10~GB, the ETL transformations (parsing and aggregation) become heavy enough to benefit from parallel CPU execution.

\subsection{Weak Scaling}
Ideal weak scaling assumes a constant runtime when resources grow at the same rate as the data. We compared H-1 (1 worker/5~GB) to W-2 (2 workers/10~GB):

\begin{itemize}
    \item \textbf{1 Worker / 5~GB (Trial W-1):} 142 seconds.
    \item \textbf{2 Workers / 10~GB (Trial W-2):} 173 seconds.
\end{itemize}

\begin{figure}[H]
    \centering
    \includegraphics[width=0.65\textwidth]{weak_scaling_bar.png}
    \caption{Weak Scaling: Maintaining 5~GB of data per worker node.}
    \label{fig:weak_scaling}
\end{figure}

\textbf{Analysis:} Weak scaling failed to maintain constant time, showing $\mathtt{\sim}$21\% increase in runtime. This degradation highlights the significant ``coordination tax'' required to manage larger datasets across multiple nodes, primarily due to increased network latency during Spark's shuffle phase and non-linear I/O performance in the virtualized HDFS environment.

\subsection{Bottleneck Identification}
A critical observation is the ``Performance Floor'' at the $\sim$80-second mark. Trials H-3 (81s), V-2 (80s), and V-4 (80s) all converged here, regardless of data size or CPU cores. Some of the possible causes are:

\begin{enumerate}
    \item \textbf{Scheduling Overhead:} At 80s, the time required to partition data and schedule tasks dominates the execution time.
    \item \textbf{Shuffle I/O:} Network bandwidth saturation during the reduce/group-by phase limits gains from faster CPUs.
    \item \textbf{Storage Limits:} Shared HDFS throughput likely hit a ceiling, capping the read speeds for both 5~GB and 10~GB datasets.
    \item \textbf{Memory Pressure:} For 10~GB loads, the RAM limit per node increases garbage collection (GC) frequency during data shuffling, further impeding scalability.
\end{enumerate}

\subsection{Speedup \& Efficiency Analysis}
Comparing our findings to theoretical linear scaling reveals a significant efficiency drop-off:

\begin{itemize}
    \item \textbf{2-Worker Case ($N=2$):} Speedup $S \approx 1.60$ (Ideal: 2.0); Efficiency $E = 80\%$.
    \item \textbf{3-Worker Case ($N=3$):} Speedup $S \approx 1.75$ (Ideal: 3.0); Efficiency $E \approx 58\%$.
\end{itemize}

The efficiency plummet at three workers confirms that the 5~GB dataset is too small to offset the overhead of the third node. According to Amdahl’s Law, the sequential portions of the code and the driver-side result collection create a fixed cost that prevents the system from achieving linear speedup regardless of additional worker allocation.